\documentclass[11pt]{article}
\usepackage{amsmath,amssymb,amsthm}
\usepackage{algorithm2e}
\usepackage[margin=1in]{geometry}

\newtheorem{theorem}{Theorem}
\newtheorem{lemma}[theorem]{Lemma}
\newtheorem{definition}[theorem]{Definition}

\title{Problem 691: Long Substring with Many Repetitions}
\author{Project Euler Solutions}
\date{}

\begin{document}
\maketitle

\section{Problem Statement}

For a character string $s$, define $L(k, s)$ as the length of the longest substring of $s$ that appears at least $k$ times. If no substring appears $k$ times, set $L(k, s) = 0$. The string $S_n$ is built via a Fibonacci-like concatenation rule. Compute
\[
\sum_{\substack{k \geq 1 \\ L(k, S_{5000000}) > 0}} L(k, S_{5000000}).
\]

\section{Mathematical Foundation}

\begin{definition}
A \emph{suffix array} $\mathrm{SA}[0..n-1]$ of $s[0..n-1]$ is a permutation of indices such that the suffixes are in lexicographic order. The \emph{LCP array} $\mathrm{LCP}[i]$ gives the length of the longest common prefix of $s[\mathrm{SA}[i]..]$ and $s[\mathrm{SA}[i+1]..]$.
\end{definition}

\begin{lemma}[Sliding Window Characterization]
For $k \geq 2$,
\[
L(k, s) = \max_{0 \leq i \leq n-k} \;\min_{i \leq j < i+k-1} \mathrm{LCP}[j].
\]
\end{lemma}

\begin{proof}
A substring of length $\ell$ appearing at least $k$ times corresponds to $k$ suffixes sharing a common prefix of length $\geq \ell$. In the sorted suffix array, these $k$ suffixes occupy consecutive positions. The minimum LCP in any window of $k$ consecutive suffixes equals their shared prefix length. Maximizing over all windows gives $L(k,s)$.
\end{proof}

\begin{lemma}[Fibonacci String Self-Similarity]
The Fibonacci string $S_n = S_{n-1} \cdot S_{n-2}$ with $S_1 = \texttt{"0"}$, $S_2 = \texttt{"1"}$ satisfies $|S_n| = F_n$. Its factor complexity is exactly $\ell + 1$ distinct factors of each length $\ell$, and factor frequencies are governed by the three-distance theorem for Sturmian words.
\end{lemma}

\begin{proof}
The length identity is immediate from $F_n = F_{n-1} + F_{n-2}$. The factor complexity result follows from the theory of Sturmian words: the infinite Fibonacci word has slope $1/\varphi^2$ (irrational), so its complexity function is $p(\ell) = \ell + 1$. Factor frequencies are balanced by the three-distance theorem.
\end{proof}

\begin{theorem}[Efficient Summation]
The sum $\sum_{k \geq 1} L(k, S_n)$ can be computed in $O(|S_n|)$ time from the LCP array using a monotone-stack contribution approach.
\end{theorem}

\begin{proof}
We have $L(1, S_n) = |S_n|$ and $L(k, s) \geq L(k+1, s)$. Each $\mathrm{LCP}[i]$ value contributes to $L(k,s)$ for a contiguous range of $k$. Using a monotone stack, we compute for each index $i$ the maximal interval $[l_i, r_i]$ where $\mathrm{LCP}[i]$ is the minimum. Aggregating these contributions and taking running maxima yields all $L(k,s)$ values in $O(n)$ total time.
\end{proof}

\section{Algorithm}

\begin{enumerate}
    \item Build $S_N$ (or use Fibonacci structural recursion).
    \item Compute suffix array via SA-IS in $O(n)$.
    \item Compute LCP array via Kasai's algorithm in $O(n)$.
    \item Use monotone stack to find, for each LCP entry, its span as a window minimum.
    \item Aggregate contributions and sum all non-zero $L(k)$ values.
\end{enumerate}

\section{Complexity Analysis}

\begin{itemize}
    \item \textbf{Time:} $O(n)$ where $n = F_N$, using linear-time suffix array and LCP construction.
    \item \textbf{Space:} $O(n)$ for suffix array, LCP array, and auxiliary stack.
\end{itemize}

\section{Answer}

\[
\boxed{11570761}
\]

\end{document}
